\chapter{Conclusions and outlook}\label{ch:conclusions}

\section{Conclusions}

\begin{itemize}
  \item Cross-track baselines turn topography into a channel-dependent phase that the standard multichannel
        reconstruction cannot absorb, $\varphi_i\approx 2\pi\bxt_{,i}h/(\lambda r_0\tan\theta)$. Without correction the median highest ambiguity of the sweeps is \SI{\SweepOneMedNone}{dB} in 1-D and
        \SI{\SweepTwoMedNone}{dB} in 2-D, against a reference near \SI{-48}{dB}.
  \item SATA computes the right phase from the DEM, but applied to each aliased channel before the reconstruction it can give
        only one phase to the parcels of positions $X\approx\SI{6.2}{km}$ apart that share a Doppler bin. It is exact only
        when the terrain has the same height at all those positions; elsewhere it leaves a structural residual that is predicted
        from the geometry within about \SI{1}{dB}, and with a DEM only it improves the median by \SI{5.5}{dB} at best.
  \item Placing the DEM phase inside the inversion solves this. The joint reconstruction over the whole line models every
        channel from the focused line and the DEM and inverts it by regularised least squares. Using only the system parameters
        and the DEM it reaches a median of \SI{\SweepOneMedJoint}{dB} in 1-D (\SweepOneWithinTen\,\% of the runs within
        \SI{10}{dB} of the reference, about \SI{90}{\percent} for four or more channels) and \SI{\SweepTwoMedJoint}{dB} in 2-D,
        where with three or more channels it is within \SI{1}{dB} of the reference in the median.
  \item The remaining failures are formations whose baselines do not separate the folded parcels (two channels, equally
        spaced baselines). The regularisation leaves the ambiguity there instead of amplifying noise.
  \item The accuracy is limited by the DEM: the phase error grows as $\bxt\,\delta h$, so long baselines need accurate DEMs.
\end{itemize}

\section{Outlook}

\begin{itemize}
  \item Validation with realistic antenna patterns, orbits and baselines that vary along the orbit, and eventually with real
        multistatic data.
  \item 2-D: distributed scenes, layover and shadow, finer range resolution (envelope interpolation) and wide swaths
        (range-dependent kernels in blocks).
  \item Joint estimation of a DEM correction, using the unexplained part of the channels as the cost.
  \item Use of the conditioning analysis (\cref{sec:mismatch}) as a design rule for the cross-track baselines of a formation.
  \item Computational cost: block-wise or GPU implementation of the model and its adjoint for full scenes.
\end{itemize}
